\documentclass{article}
\usepackage[T2A]{fontenc}
\usepackage[utf8]{inputenc}
\usepackage{amsthm}
\usepackage{amsmath}
\usepackage{amssymb}
\usepackage{amsfonts}
\usepackage{mathrsfs}
\usepackage[12pt]{extsizes}
\usepackage{fancyvrb}
\usepackage{indentfirst}
\usepackage[
  left=2cm, right=2cm, top=2cm, bottom=2cm, headsep=0.2cm, footskip=0.6cm, bindingoffset=0cm
]{geometry}
\usepackage[russian]{babel}


\begin{document}
\section*{Вариант 8}
Решение уравнения имеет вид

\begin{equation}
  \mathbf{U}_{\xi0} = \frac{1}{2\varepsilon^2}\left[ \frac{\partial^2 P_0}{\partial\xi\partial\tau} + 
  \frac{\partial^2 P_0}{\partial\xi\partial\tau} \bar{\Psi}(\zeta) + 
  \frac{\partial P_0}{\partial\xi}\bar{\Phi}(\zeta)  \right] ,
\end{equation}
где
\begin{align*}
  \bar{\Psi}(\zeta) = F_2(\varepsilon\zeta)D_1 - F_1(\varepsilon\zeta) - 2F_4(\varepsilon\zeta)D_2, \quad
  \bar{\Phi}(\zeta) =  2F_3(\varepsilon\zeta) - F_2(\varepsilon\zeta)D_2 - 2F_4(\varepsilon\zeta)D_1,
\end{align*}
\begin{align*}
  \bar{\Psi_1}(\zeta) = \int\limits_{\zeta}^{1} \bar{\Psi}d\zeta,  \quad
  \bar{\Phi_1}(\zeta) = \int\limits_{\zeta}^{1}\bar{\Phi}(\zeta)d\zeta, \quad   
  \varepsilon(\omega) = \sqrt{\mathfrak{R}/2}, 
\end{align*}
\end{document}

